\section{Introducción}

\subsection{Planteo del problema}

El análisis de series temporales persigue dos objetivos que conviene distinguir. El primero es
descriptivo: comprender qué estructura gobierna la evolución de una magnitud a lo largo del
tiempo, si esa estructura es estable y qué componentes la integran. El segundo es predictivo:
anticipar valores futuros con una medida explícita de la incertidumbre asociada. Este trabajo
aborda ambos sobre tres series de naturaleza distinta.

La elección de series heterogéneas responde a un propósito deliberado. Un conjunto de datos
que exhibe estacionalidad marcada y otro que carece por completo de ella exigen caminos de
identificación diferentes, y contrastarlos ilumina mejor el método que repetir tres veces el
mismo procedimiento sobre series equivalentes.

\subsection{Origen de los datos}

\subsubsection{Series de infraestructura productiva}

Las dos primeras series provienen de la telemetría operativa de una plataforma comercial en
producción y fueron registradas durante 2026.

La primera contabiliza las solicitudes atendidas por hora por el balanceador de carga que
recibe todo el tráfico entrante de la plataforma. Abarca las 744 horas del mes de julio de
2026, expresadas en tiempo universal coordinado, sin interrupciones ni huecos. El volumen promedio ronda los 2,1 millones de solicitudes por hora
\parencite{datos_alb}.

La segunda contabiliza las solicitudes por hora dirigidas a cada destino de la interfaz de
punto de venta. Abarca 720 horas comprendidas entre el 18 de julio y el 17 de agosto de 2026,
también sin faltantes, con un volumen promedio cercano a las 830.000 solicitudes por destino
y por hora \parencite{datos_pos}.

Ambas fueron reindexadas sobre una grilla horaria completa antes del análisis. El
procedimiento cumple una función diagnóstica además de correctiva: revela cualquier hueco
temporal que de otro modo pasaría inadvertido, puesto que los modelos suponen observaciones
equiespaciadas y una interrupción no declarada distorsionaría las funciones de autocorrelación.

\subsubsection{Serie de telemetría de vuelo}

La tercera serie procede del acervo ALFA, un conjunto de registros de telemetría de una
aeronave no tripulada de ala fija recopilado por el Instituto de Robótica de la Universidad
Carnegie Mellon \parencite{keipour2020}. El protocolo de comunicación transmite el estado de la
aeronave a la estación de tierra y almacena cada mensaje precedido de una marca temporal con
resolución de microsegundos.

El acervo fue construido con una finalidad específica: proveer datos para la investigación en
\textbf{detección de fallas y anomalías} en vehículos aéreos no tripulados. Comprende cuarenta y
siete vuelos, veintitrés de los cuales presentan falla súbita de motor y veinticuatro
corresponden a fallas de superficies de control \parencite{keipour2021}. Esa finalidad refuerza el interés analítico de la altitud: su predicción de corto plazo es justamente el insumo sobre el cual un sistema de supervisión detecta desviaciones respecto del comportamiento
esperado.

Los registros de telemetría en bruto, a diferencia de las secuencias procesadas del acervo, no
incorporan etiquetas de falla, de modo que no cabe afirmar a partir de ellos si un vuelo
determinado corresponde a un escenario nominal o a uno con falla inyectada. El vuelo analizado
en este trabajo no exhibe, en todo caso, indicios de falla de motor: el acelerador permanece
activo durante todo el tramo, con una media que oscila entre el veinte y el cuarenta y cinco por
ciento según el minuto, y la velocidad aerodinámica media de 16,6 metros por segundo se mantiene
próxima a la velocidad de crucero de quince metros por segundo que fijan los parámetros de
configuración. La colección disponible comprende cuarenta y un registros repartidos en cinco
jornadas de 2018.

La variable analizada es la altitud relativa al punto de despegue, transmitida en milímetros
y convertida a metros. Se prefirió esta magnitud sobre la altitud respecto del nivel del mar
por tres razones: se lee directamente como altura de vuelo, se muestrea a mayor
frecuencia, y permite detectar el contacto con el suelo sin conocer la elevación del campo.

La construcción de la serie exigió dos decisiones. La primera consistió en aislar el tramo
efectivamente en vuelo, puesto que la mayor parte de cada registro corresponde a la aeronave
detenida en tierra, donde la altitud transmitida es únicamente deriva del barómetro. En uno de
los vuelos examinados, el registro se extiende durante 2.880 segundos mientras la mediana de
velocidad respecto del suelo alcanza apenas 0,16 metros por segundo. Modelar ese tramo
equivaldría a modelar ruido de sensor. Se conservó, en consecuencia, el bloque contiguo más
extenso con altitud superior a tres metros.

La segunda decisión fue la frecuencia de análisis. La serie se llevó a una grilla regular de
un segundo. Dado que la frecuencia nativa ronda los 4,15 hertz, la operación es un submuestreo y no introduce puntos interpolados. Una grilla más fina habría duplicado el número
de observaciones al costo de interpolar, y esos puntos artificiales inflarían la
autocorrelación de rezago corto, que es precisamente la magnitud sobre la cual se apoya la
identificación del modelo.

La serie resultante comprende 457 observaciones que recorren un perfil de vuelo completo,
desde el despegue hasta el descenso, con altitudes entre 4,9 y 82,8 metros.

\subsection{Organización del informe}

El apartado siguiente expone el marco teórico. El posterior desarrolla el análisis, organizado
según los doce puntos que estructuran el trabajo. Las conclusiones sintetizan los hallazgos y
delimitan su alcance. Los apéndices reúnen el material gráfico complementario y las rutinas de
cómputo extensas.
